% GENERATED FILE. Edit canonical lessons, then run npm run generate:books.
% canonical-source-hash: fnv1a64:86603e82eee0d338
% canonical-lessons: TA-C37-uur, TA-C37-nanban, TA-C37-ivar, TA-C37-ivar-en-nanbar

\chapter{Your Town, Your Friend}
\label{ch:ta-your-town-your-friend}
\hlchaptermodality{37}

\begin{chapteropening}
\textbf{By the end of this chapter:} \emph{I can ask \ta{உங்கள்} \ta{ஊர்} \ta{என்ன}? and introduce someone with \ta{இவர்} \ta{என்} \ta{நண்பர்}, choose between \ta{நண்பன்}, \ta{நண்பி} and \ta{நண்பர்}, run Tamil's \ta{இ} / \ta{அ} / \ta{எ} pointing series, and hear \ta{ண்} voice the \ta{ப} that follows it.}

Say \ta{இவர்} \ta{என்} \ta{நண்பர்} with no verb in it, pick the ending by respect rather than by habit, name the -ūr inside Tanjore and Bengaḷūru, and state what Tamil forces you to decide for a sibling that it does not force for a friend.
\end{chapteropening}

\section[ūr]{\ta{ஊர்} (ūr) — \textquotedblleft{}town,\textquotedblright{} and the question Tamil asks early}

\label{lesson:TA-C37-uur}



\begin{quote}
\emph{Eṉ tavaṟu} — my mistake, a word that is a verb and a noun at once. Then the general word for food. (\emph{Uṇavu}.) Now back to meeting somebody.
\end{quote}

\subsection*{You'll want to know: \ta{ஊர்}}
\begin{quote}
\textbf{\ta{ஊர்}} — \emph{ūr} — \textbf{town, village; the place you are from}
\end{quote}

That last sense matters most. \textbf{\ta{உங்கள்} \ta{ஊர்}} is your \textbf{native place}, where your family is from, which may be somewhere you have never lived.

\begin{quote}
\textbf{\ta{உங்கள்} \ta{ஊர்} \ta{என்ன}} — \emph{uṅgaḷ ūr eṉṉa} — \textbf{What is your town}
\end{quote}

It lands early in a Tamil conversation and is not nosiness — the slot English fills with "what do you do.

\subsection*{The letters in this word}
\textbf{\ta{ஊ}} is a standing vowel — long \emph{ū}, partner of the \textbf{\ta{உ}} that opens \emph{uṇavu}. \textbf{\ta{ர்}} is \textbf{\ta{ர}}, the soft r, shut by a puḷḷi. Hold the vowel: \emph{ūr}, not \emph{ur}. Vowel length is part of a word's identity here.

\begin{grammarlens}[title={Grammar Lens: the two registers, side by side}]
You met this choice on the age question:

\noindent
\begin{tabularx}{\linewidth}{@{}>{\raggedright\arraybackslash}X>{\raggedright\arraybackslash}X@{}}
\toprule
\textbf{the question} & \textbf{who you say it to} \\
\midrule
\textbf{\ta{உன்} \ta{ஊர்} \ta{என்ன}} \emph{uṉ ūr eṉṉa} & a friend, a child, a peer \\
\textbf{\ta{உங்கள்} \ta{ஊர்} \ta{என்ன}} \emph{uṅgaḷ ūr eṉṉa} & a stranger, an elder, anyone senior \\
\bottomrule
\end{tabularx}

Both are possessive plus noun plus question word, with \textbf{no verb} — the shape you learned for \textbf{\ta{உன்} \ta{வயசு} \ta{என்ன}}. Only the possessive changes, and with it the social temperature. When unsure, use \textbf{\ta{உங்கள்}}.
\end{grammarlens}

\begin{cousinweb}
\textbf{\ta{ஊர்}} is inherited Dravidian, and welded into city names you know in English:

\noindent
\begin{tabularx}{\linewidth}{@{}>{\raggedright\arraybackslash}X>{\raggedright\arraybackslash}X@{}}
\toprule
\textbf{the English name} & \textbf{what it is} \\
\midrule
\textbf{Tanjore} & Tamil \textbf{\ta{தஞ்சாவூர்}} \emph{Tañjāvūr} \\
\textbf{Vellore} & Tamil \textbf{\ta{வேலூர்}} \emph{Vēlūr} \\
\textbf{Bangalore} & Kannada \textbf{\ka{ಬೆಂಗಳೂರು}} \emph{Bengaḷūru} \\
\bottomrule
\end{tabularx}

Every one of those \emph{-ore} endings is this word: the British ear turned \emph{-ūr} into \emph{-ore}, and the cities have taken their own names back.

Two exactnesses. Bangalore is a \textbf{Kannada} town, so its \emph{-ūru} is Kannada's form of the inherited word — the family's, not Tamil's alone. And the \emph{first} halves are not all settled; \emph{Vēl-} in Vellore is argued over. The ending is the certain part.

Weigh that against the rice word. \textbf{\ta{அரிசி}} \emph{may} stand behind English \emph{rice} through Greek, and the dictionaries hedge. Nobody hedges about \emph{-ūr}.
\end{cousinweb}

\subsection*{Your turn}
Say these aloud:

\begin{itemize}
\raggedright
  \item the long vowel — \textquotedblleft{}ūr,\textquotedblright{} not \textquotedblleft{}ur\textquotedblright{}
  \item to a stranger — \textquotedblleft{}uṅgaḷ ūr eṉṉa\textquotedblright{}
  \item to a friend — \textquotedblleft{}uṉ ūr eṉṉa\textquotedblright{}
  \item the ending inside the map — \textquotedblleft{}Tañjāvūr … Vēlūr … Bengaḷūru\textquotedblright{}
\end{itemize}

\begin{tcolorbox}[breakable,colback=teal!4,colframe=teal!35!black,title={Before you move on}]
Ask a stranger where they are from, then a friend, and say what \textbf{\ta{ஊர்}} means beyond \textquotedblleft{}town.\textquotedblright{} (\emph{Uṅgaḷ ūr eṉṉa}; \emph{uṉ ūr eṉṉa}; \textbf{your native place}.) Which earlier question has this shape, and is \emph{-ore} in Tanjore really this word? (\textbf{\ta{உன்} \ta{வயசு} \ta{என்ன}}; \textbf{yes}.) How sure are we there, next to \emph{arisi} and \emph{rice}? (\textbf{Much surer.})
\end{tcolorbox}
\section[naṇbaṉ]{\ta{நண்பன்} (naṇbaṉ) — \textquotedblleft{}friend,\textquotedblright{} with the ending doing the work}

\label{lesson:TA-C37-nanban}



\begin{quote}
A stranger's native place, then the ending inside Tanjore. (\emph{Uṅgaḷ ūr eṉṉa}; \emph{-ūr}.) Now the person you would introduce.
\end{quote}

\subsection*{You'll want to know: \ta{நண்பன்}}
\begin{quote}
\textbf{\ta{நண்பன்}} — \emph{naṇbaṉ} — \textbf{friend} (a man)
\end{quote}

Tamil will not let you say a bare \textquotedblleft{}friend\textquotedblright{} any more than a bare \textquotedblleft{}brother.\textquotedblright{} The stem is \textbf{\ta{நண்ப}-}; the ending decides.

\noindent
\begin{tabularx}{\linewidth}{@{}>{\raggedright\arraybackslash}X>{\raggedright\arraybackslash}X@{}}
\toprule
\textbf{the form} & \textbf{who it is} \\
\midrule
\textbf{\ta{நண்பன்}} \emph{naṇbaṉ} & a male friend \\
\textbf{\ta{நண்பி}} \emph{naṇbi} & a female friend \\
\textbf{\ta{நண்பர்}} \emph{naṇbar} & a friend spoken of with respect \\
\bottomrule
\end{tabularx}

The \textbf{-\ta{அர்}} is the respectful plural that turns \textbf{\ta{நீ}} into \textbf{\ta{நீங்கள்}}, applied to a person.

This is a different axis from the family words. \textbf{\ta{அண்ணன்}} and \textbf{\ta{தம்பி}} split brothers by \textbf{age}; the friend words split by \textbf{sex}, \textbf{\ta{நண்பர்}} by \textbf{respect}.

\subsection*{The letters in this word}
\textbf{\ta{ந}} + \textbf{\ta{ண்}} + \textbf{\ta{ப}} + \textbf{\ta{ன்}} — \textbf{every one of Tamil's n-letters inside one word}. \textbf{\ta{ந}} opens it, dental, at the teeth. \textbf{\ta{ண்}} is retroflex, tongue curled back, closed by a puḷḷi. \textbf{\ta{ன்}} ends it, the alveolar n from \textbf{\ta{நன்றி}}. The writing lesson gave them to you one at a time; here they arrive together.

\begin{sounds}
The letter after \textbf{\ta{ண்}} is \textbf{\ta{ப}}, named \emph{pa}. This book writes \emph{naṇ\textbf{b}aṉ} because of the \textbf{\ta{நன்றி}} rule: a nasal in front of a stop \textbf{voices it}. A Tamil reader supplies the voicing, as an English reader supplies the \emph{sh} in \emph{sugar}.
\end{sounds}

\begin{cousinweb}
Behind \textbf{\ta{நண்பன்}} is \textbf{\ta{நண்பு}} (\emph{naṇbu}), \textquotedblleft{}friendship, closeness,\textquotedblright{} which belongs to the root family of \textbf{\ta{நண்ணு}} (\emph{naṇṇu}), \textquotedblleft{}\textbf{to draw near}.\textquotedblright{} A friend, in Tamil's own accounting, is someone who came close.

Be precise. The Dravidian dictionary files \emph{naṇbu} and \emph{naṇbaṉ} in the same entry as \emph{naṇṇu}; the Tamil lexicon prefers \textbf{\ta{நள்}-}, making them siblings rather than parent and child. They differ on the line of descent, not on the household.

Set it beside \textbf{\ta{அப்பா}} and \textbf{\ta{அம்மா}}, not built out of anything — the repeated-syllable shape children produce for parents worldwide. \textbf{\ta{நண்பன்}} is assembled, from a root you can point at.
\end{cousinweb}

\subsection*{Your turn}
Say these aloud:

\begin{itemize}
\raggedright
  \item the three endings — \textquotedblleft{}naṇbaṉ, naṇbi, naṇbar\textquotedblright{}
  \item the voicing after the nasal — \textquotedblleft{}naṉṟi … naṇbaṉ\textquotedblright{}
  \item the root it came from — \textquotedblleft{}naṇṇu, to draw near\textquotedblright{}
  \item the family words that split on age instead — \textquotedblleft{}aṇṇaṉ, tambi\textquotedblright{}
\end{itemize}

\begin{tcolorbox}[breakable,colback=teal!4,colframe=teal!35!black,title={Before you move on}]
Give the friend-word for a woman, a man, and someone you respect. (\emph{Naṇbi}; \emph{naṇbaṉ}; \emph{naṇbar}.) Which axis do the friend words split on, and which do the brother words? (\textbf{Sex and respect}; \textbf{age}.) Why is \textbf{\ta{ப}} heard as \emph{b}, and are \textbf{\ta{அப்பா}} and \textbf{\ta{அம்மா}} built from a root the way \emph{naṇbaṉ} is? (\textbf{A nasal voices the stop after it}; \textbf{no}.)
\end{tcolorbox}
\section[ivar]{\ta{இவர்} (ivar) — \textquotedblleft{}this person,\textquotedblright{} and Tamil's pointing system}

\label{lesson:TA-C37-ivar}



\begin{quote}
\emph{Naṇbaṉ}, \emph{naṇbi}, \emph{naṇbar} — and the root behind them. (\emph{Naṇṇu}, to draw near.) You have the word for a friend, but cannot point at one.
\end{quote}

\subsection*{You'll want to know: \ta{இவர்}}
\begin{quote}
\textbf{\ta{இவர்}} — \emph{ivar} — \textbf{this person} (near me, spoken of respectfully)
\end{quote}

Not \textquotedblleft{}he\textquotedblright{} and not \textquotedblleft{}she\textquotedblright{}: \textbf{\ta{இவர்}} is what you say about a person standing here whom you are treating with respect, and it does not commit you to their sex. It is built on \textbf{\ta{இவன்}} (\emph{ivaṉ}) and \textbf{\ta{இவள்}} (\emph{ivaḷ}), with the same \textbf{-\ta{அர்}} that gave you \textbf{\ta{நண்பர்}}.

One caution about speech: \textbf{\ta{இவர்}} is heard mostly of men, and \textbf{\ta{இவங்க}} (\emph{ivaṅga}) does the same respectful job for a woman.

\subsection*{The letters in this word}
\textbf{\ta{இ}} + \textbf{\ta{வ}} + \textbf{\ta{ர்}}. \textbf{\ta{இ}} is the standing vowel \emph{i} — the letter, not the \textbf{\ta{ி}} sign that hangs on a consonant. \textbf{\ta{வ}} is the \emph{va} opening \textbf{\ta{வணக்கம்}} and \textbf{\ta{வீடு}}. \textbf{\ta{ர்}} is the soft r with a puḷḷi.

A word this short is almost all arc, which is the moment to recall the usual explanation for that roundness: Tamil was scratched into palm leaf, and a straight stroke along the grain can split the leaf. Hold it as the writing lesson asked — \textbf{the standard account, not a settled fact}, since the earliest Tamil writing was cut into stone and is angular. What survives is the useful part: \textbf{the tool leaves fingerprints on the letters}.

\begin{grammarlens}[title={Grammar Lens: near, far, and which}]
\textbf{\ta{இவர்}} is one cell of a system, and the system is a single letter at the front:

\noindent
\begin{tabularx}{\linewidth}{@{}>{\raggedright\arraybackslash}X>{\raggedright\arraybackslash}X@{}}
\toprule
\textbf{the front letter} & \textbf{what it does} \\
\midrule
\textbf{\ta{இ}-} & points \textbf{near} — this one, here \\
\textbf{\ta{அ}-} & points \textbf{away} — that one, there \\
\textbf{\ta{எ}-} & \textbf{asks} — which one \\
\bottomrule
\end{tabularx}

So \textbf{\ta{இவர்}} is \textquotedblleft{}this person,\textquotedblright{} \textbf{\ta{அவர்}} (\emph{avar}) \textquotedblleft{}that person,\textquotedblright{} \textbf{\ta{எவர்}} (\emph{evar}) \textquotedblleft{}which person.\textquotedblright{} You have been using one front already: \textbf{\ta{என்ன}} (\emph{eṉṉa}), \textquotedblleft{}what,\textquotedblright{} is the \textbf{\ta{எ}-} of asking.

Two honest notes. \textbf{\ta{எ}-} is the \textbf{question}, not a third distance. And Old Tamil had a genuine third distance, \textbf{\ta{உ}-}, for something visible but far off; modern Tamil has retired it.
\end{grammarlens}

\subsection*{Your turn}
Say these aloud:

\begin{itemize}
\raggedright
  \item the three fronts — \textquotedblleft{}i- … a- … e-\textquotedblright{}
  \item the row they build — \textquotedblleft{}ivar, avar, evar\textquotedblright{}
  \item the question word you already had — \textquotedblleft{}eṉṉa\textquotedblright{}
  \item the respectful ending on both — \textquotedblleft{}naṇbar … ivar\textquotedblright{}
\end{itemize}

\begin{tcolorbox}[breakable,colback=teal!4,colframe=teal!35!black,title={Before you move on}]
What do \textbf{\ta{இ}-}, \textbf{\ta{அ}-} and \textbf{\ta{எ}-} do at the front of a word, and which question word you already knew is built on \textbf{\ta{எ}-}? (\textbf{Near}; \textbf{away}; \textbf{asks} — and \textbf{\ta{என்ன}}.) Does \textbf{\ta{இவர்}} tell you whether the person is a man or a woman? (\textbf{No} — it marks \textbf{respect}; spoken Tamil leans \textbf{\ta{இவங்க}} for a woman.) How firmly should you hold the palm-leaf explanation for the curves? (As \textbf{the standard account, not a settled fact}.)
\end{tcolorbox}
\section[ivar eṉ naṇbar]{\ta{இவர்} \ta{என்} \ta{நண்பர்} — \textquotedblleft{}this is my friend\textquotedblright{}}

\label{lesson:TA-C37-ivar-en-nanbar}



\begin{quote}
\emph{Ivar, avar, evar} — near, away, the question. Then \emph{naṇbaṉ, naṇbi, naṇbar}. Everything this lesson needs is in your mouth; what is left is the order.
\end{quote}

\subsection*{You'll want to know: the introduction}
\begin{quote}
\textbf{\ta{இவர்} \ta{என்} \ta{நண்பர்}.} — \emph{ivar eṉ naṇbar} — \textbf{This is my friend.}
\end{quote}

Three words, and \textbf{not one of them is a verb}. Tamil links two nouns by setting them beside each other, as in \textbf{\ta{என்} \ta{பெயர்} …} and \textbf{\ta{என்} \ta{தவறு}}.

The pieces in order: \textbf{\ta{இவர்}}, the person here; \textbf{\ta{என்}}, my; \textbf{\ta{நண்பர்}}, friend-with-respect. Then stop.

The whole meeting:

\begin{quote}
— \textbf{\ta{வணக்கம்}. \ta{இவர்} \ta{என்} \ta{நண்பர்}.}
\end{quote}

>

\begin{quote}
— \textbf{\ta{வணக்கம்}. \ta{உங்கள்} \ta{ஊர்} \ta{என்ன}.}
\end{quote}

Every word in it is yours.

\subsection*{The letters in this word}
\textbf{\ta{என்}} is \textbf{\ta{எ}} with \textbf{\ta{ன்}}, the alveolar n dotted shut. \textbf{\ta{நண்பர்}} opens on dental \textbf{\ta{ந}}, turns on retroflex \textbf{\ta{ண்}}, closes on \textbf{\ta{ர்}}. So one line makes you produce \textbf{\ta{ன்}} and \textbf{\ta{ண்}} within a breath of each other — one at the ridge, the other with the tongue curled back. Hear that difference and the three-way n is no longer theory.

\begin{culture}
\textbf{\ta{நண்பர்}} rather than \textbf{\ta{நண்பன்}} is a decision, not a default.

\begin{itemize}
\raggedright
  \item \textbf{\ta{இவர்} \ta{என்} \ta{நண்பன்}.} — a male friend, plainly stated; warm among peers.
  \item \textbf{\ta{இவர்} \ta{என்} \ta{நண்பி}.} — a female friend, the matching plain form.
  \item \textbf{\ta{இவர்} \ta{என்} \ta{நண்பர்}.} — the respectful ending; safe with anyone, and right when either party is older or senior.
\end{itemize}

If unsure, take \textbf{\ta{நண்பர்}}. Over-respect is invisible; under-respect is not.

Notice how differently Tamil makes you think from one word to the next. For a \textbf{sibling} it forces \textbf{age}: there is no bare \textquotedblleft{}brother\textquotedblright{} — you must choose \textbf{\ta{அண்ணன்}} or \textbf{\ta{தம்பி}}, \textbf{\ta{அக்கா}} or \textbf{\ta{தங்கை}}, and say whether they are older or younger than you before you can say anything at all. For a \textbf{friend} it forces \textbf{sex or respect}. For \textbf{parents}, nothing: \textbf{\ta{அப்பா}} and \textbf{\ta{அம்மா}} came into the language whole. Tamil rarely lets you be vague about a relationship.
\end{culture}

\subsection*{Your turn}
Say these aloud:

\begin{itemize}
\raggedright
  \item the introduction — \textquotedblleft{}ivar eṉ naṇbar\textquotedblright{}
  \item the same line to a peer about a male friend — \textquotedblleft{}ivar eṉ naṇbaṉ\textquotedblright{}
  \item the greeting and the question around it — \textquotedblleft{}vaṇakkam … uṅgaḷ ūr eṉṉa\textquotedblright{}
  \item what a sibling word forces you to decide — \textquotedblleft{}aṇṇaṉ, older; tambi, younger\textquotedblright{}
\end{itemize}

\begin{tcolorbox}[breakable,colback=teal!4,colframe=teal!35!black,title={Before you move on}]
Introduce a friend respectfully, and say how many verbs the sentence has. (\emph{Ivar eṉ naṇbar}; \textbf{none}.) What does a sibling word force you to decide that a friend word does not? (\textbf{Whether they are older or younger} — \emph{aṇṇaṉ} against \emph{tambi}.) Which friend ending is safe when unsure, and what does \textbf{\ta{இ}-} tell your listener? (\textbf{\ta{நண்பர்}}; that the person is \textbf{near}.)
\end{tcolorbox}
